\section{Conclusions and Future Work}

We have shown how refinement reflection---namely
reflecting the definitions of functions in their
output refinements---can be used to convert a
legacy programming language, like Haskell,
into a theorem prover.
%
Reflection ensures that (refinement) type checking
stays decidable and predictable via careful design
of the logic and proof combinators.
%
Reflection enables programmers working with
the highly tuned libraries, compilers and
run-times of legacy languages to specify
and verify arbitrary properties of their
code simply by writing programs in the
legacy language.
%
Our evaluation shows that refinement reflection
lets us prove deep specifications of a variety
of implementations, for example, the determinism
of fast parallel programming libraries, and also
identifies two important avenues for research.

First, while our proofs are often elegant
and readable, they can sometimes be cumbersome.
For example, in the proof of associativity of
the monadic bind operator for the @Reader@ monad
three out of eight (extensional) equalities required
explanations, some nested under multiple
$\lambda$-abstractions. Thus, it would be
valuable to explore how to extend the notions
of tactics, proof search, and automation
to the setting of legacy languages.
%
Similarly, while our approach to $\alpha$- and
$\beta$-equivalence is sound, we do not know
if it is \emph{complete}. We conjecture it is,
due to the fact that our refinement terms
are from the simply typed lambda calculus (STLC).
Thus, it would be interesting to use the
normalization of STLC to develop a sound
and complete system for SMT-based type-level
computation and use it to automate proofs
predictably.

Second, while Haskell's separation of pure
and effectful code undoubtedly makes it
easier to implement refinement reflection,
we believe that the technique, like
refinement typing in general, is orthogonal
to purity.
In particular, by carefully mediating
between the interaction of pure and
impure code using methods like
permissions, uniqueness \etc
refinement type systems have been
developed for legacy languages
like C \cite{LiquidPOPL10},
JavaScript \cite{Chugh2012, Vekris16},
and Racket \cite{RefinedRacket}
and so it would be interesting
to see how to extend refinement
reflection to other legacy languages.
